\documentclass[14pt,a4paper]{extarticle}
\usepackage[T2A]{fontenc}\usepackage[utf8]{inputenc}\usepackage[russian]{babel}
\usepackage{geometry}\geometry{left=30mm,right=15mm,top=20mm,bottom=20mm}
\setlength{\parindent}{1.25cm}\setlength{\parskip}{0.3em}\sloppy
\begin{document}\begin{titlepage}\begin{center}
САНКТ-ПЕТЕРБУРГСКИЙ НАЦИОНАЛЬНЫЙ\\ ИССЛЕДОВАТЕЛЬСКИЙ УНИВЕРСИТЕТ ИТМО
\vspace{1.5cm}

Дисциплина: Бэк-энд разработка
\vspace{1.5cm}

\textbf{Отчёт}\\[0.5cm]
Домашняя работа 9\\[0.5cm]
\textbf{Документирование веб-приложения}
\vfill
\begin{flushright}Выполнил:\\ Цатинян Артём\\ Группа БР1.2\\[0.8cm] Проверил:\\ Добряков Д. И.\end{flushright}
\vfill Санкт-Петербург\\ 2026 г.\end{center}\end{titlepage}
\section{Задача и состояние}
Описать документацию микросервисного приложения с помощью OpenAPI, Postman и MkDocs. Отчёт содержит обзор имеющихся материалов и проект их объединения; публикация сайта документации отдельно не проверялась.
\section{Состав приложения}
Приложение включает сервисы identity, catalog, booking и review. В ЛР5 доступ к их HTTP-интерфейсам проверялся через nginx; в проекте имеются спецификации API, а результаты проверки описаны в отчёте ЛР5.
\section{OpenAPI и Swagger}
Для каждого сервиса следует описать адреса и методы, параметры, тела запросов и ответов, коды ошибок и авторизацию. Спецификации должны соответствовать фактическим контроллерам. Swagger UI позволяет изучать операции и выполнять запросы в доступном окружении; корректность серверного адреса и схем проверяется отдельно.
\section{Коллекция Postman}
Предлагается сгруппировать запросы по сервисам и вынести базовый адрес и токен в переменные окружения. Демонстрационный сценарий: авторизация, получение данных каталога, бронирование и просмотр результата. Для проверок указываются ожидаемые коды и поля; реальные значения токенов в коллекцию не включаются.
\section{Сайт MkDocs}
В структуру документации предлагается включить обзор архитектуры, инструкции локального запуска, разделы сервисов, взаимодействие и очереди, спецификации API, развёртывание и мониторинг. Ссылки на реальные файлы и команды следует проверить при сборке сайта.
\section{Вывод}
Подготовлена структура единой документации; сборка MkDocs, коллекция и проверка актуальности требуют отдельной реализации.
\end{document}
